\documentclass[aspectratio=169]{beamer}
\usepackage{ctex, hyperref}
%\usepackage[T1]{fontenc}
% \usefonttheme{professionalfonts} % 使用专业字体主题

% 在加载 fontspec 之前设置字体族
\renewcommand{\sfdefault}{cmss} % 设置无衬线字体族为 Computer Modern Sans Serif
\usepackage{algorithm,algorithmic} 
\usepackage{fontspec}
% other packages
\usepackage{latexsym,amsmath,xcolor,multicol,booktabs,calligra,fixdif}
\usepackage{graphicx,pstricks,listings,stackengine,threeparttable,multirow}
\usepackage{times}

\usepackage{tikz}
\usepackage{pgf}
\usepackage{pgfpages}
\usepackage{gensymb}
%\usepackage{newtxmath}
\usepackage[english]{babel}
%\usepackage[numbers,sort&compress,square,super]{natbib}
\usepackage{cite}
\usepackage{breqn}
% \usepackage{indentfirst} 
% \setlength{\parindent}{2em}
\usepackage{caption,tcolorbox,braket}  % 引入 caption 宏包
\usepackage{pifont}
\captionsetup{
    justification=centering,  % 标题居中对齐
    font=tiny,
    labelsep=colon,  % 标签和标题之间的分隔符（可选）
    indention=0pt  % 取消 caption 的缩进
}
\usepackage{enumitem}  % For customizing lists
\linespread{1.25} % 设置1.5倍行距
\usepackage{ragged2e}
% \usepackage{algorithm,algorithmic} 
\justifying\let\raggedright\justifying
\setbeamertemplate{bibliography item}[text]
\newcommand{\upcite}[1]{\textsuperscript{\textsuperscript{\cite{#1}}}}

\usetikzlibrary{tikzmark}
% \tcbuselibrary{skins, breakable, theorems}
\tcbuselibrary{most}
\newtcbtheorem[number within=section]{aim}{}%
  {enhanced, breakable,
    colback = deepblue!10, colframe = deepblue, colbacktitle = deepblue,
    attach boxed title to top left = {yshift = -2.5mm, xshift = 5mm},
    boxed title style = {sharp corners},
    fonttitle = \sffamily\bfseries, separator sign = {~}}{qst}

\author[李四]{李四}

\title{季度工作报告通用PPT模版}
\subtitle{框架完整的年终总结、工作汇报、述职报告PPT模版}
\institute[核物理基础研究室]{核物理基础研究室}
\date{\today}
\usepackage{dlut}

% defs
\def\cmd#1{\texttt{\color{red}\footnotesize $\backslash$#1}}
\def\env#1{\texttt{\color{blue}\footnotesize #1}}
\definecolor{deepblue}{RGB}{51,70,139}
\definecolor{deepred}{rgb}{0.6,0,0}
\definecolor{deepgreen}{rgb}{0,0.5,0}
\definecolor{halfgray}{gray}{0.55}

\lstset{
    basicstyle=\ttfamily\small,
    keywordstyle=\bfseries\color{deepblue},
    emphstyle=\ttfamily\color{deepred},    % Custom highlighting style
    stringstyle=\color{deepgreen},
    numbers=left,
    numberstyle=\small\color{halfgray},
    rulesepcolor=\color{red!20!green!20!blue!20},
    frame=shadowbox,
}   

% 设置英文字体为 Times New Roman
%\setmainfont{Times New Roman}

\begin{document}

\begin{frame}[plain]
    \titlepage
\end{frame}

\maketotalcontents

% ===================== Section 1 =====================
\setsectionsubtitle{Template Overview}
\section{模板介绍}

\begin{frame}{文件结构与编译方式}
\begin{columns}[T]
\begin{column}{0.50\textwidth}
\begin{aim*}{模板文件清单}{}
\begin{itemize}[left=0pt,label=\textcolor{deepblue}{\ding{237}}]
    \item \texttt{slide.tex} \enspace 主文档，在此编辑内容
    \item \texttt{DLUT.sty} \enspace 主题样式，一般无需修改
    \item \texttt{ref.bib} \enspace 参考文献数据库
    \item \texttt{logo.png} \enspace 封面右上角徽标
    \item \texttt{background1.png} \enspace 封面背景横幅
    \item \texttt{seccon.png} \enspace 章节过渡页背景
    \item \texttt{pic/logo\_ciae.png} \enspace 页眉徽标
\end{itemize}
\end{aim*}
\end{column}
\begin{column}{0.46\textwidth}
\begin{aim*}{编译方式（必须用 XeLaTeX）}{}
模板使用微软雅黑字体，\textbf{必须}用 XeLaTeX 编译。目录、TikZ 覆盖层等交叉引用需多次编译才能正确定位，\textbf{推荐}使用 \texttt{latexmk} 自动处理：
\medskip

\texttt{\small latexmk -xelatex slide.tex}
\medskip

如手动编译，需连续执行\textbf{三次}：\\[0.2em]
\texttt{\small xelatex slide.tex}\\[0.2em]
\texttt{\small xelatex slide.tex}\\[0.2em]
\texttt{\small xelatex slide.tex}
\end{aim*}
\end{column}
\end{columns}
\end{frame}

\begin{frame}{封面信息与章节副标题}
\begin{columns}[T]
\begin{column}{0.50\textwidth}
\textbf{封面信息} — 修改 \texttt{slide.tex} 头部的以下命令：
\medskip
\begin{tabular}{ll}
\toprule
\textbf{命令} & \textbf{作用} \\
\midrule
\texttt{\textbackslash author[简称]\{全名\}} & 汇报人姓名 \\
\texttt{\textbackslash title\{...\}}         & 报告主标题 \\
\texttt{\textbackslash subtitle\{...\}}      & 报告副标题 \\
\texttt{\textbackslash institute[简称]\{全称\}} & 单位名称 \\
\texttt{\textbackslash date\{\textbackslash today\}} & 日期 \\
\bottomrule
\end{tabular}
\end{column}
\begin{column}{0.46\textwidth}
\begin{aim*}{章节副标题用法}{}
在每个 \cmd{section} \textbf{前一行}调用：
\smallskip

{\small\texttt{\cmd{setsectionsubtitle}\{副标题\}}}\\[0.2em]
{\small\texttt{\cmd{section}\{章节名称\}}}
\smallskip

副标题显示在章节过渡页标题下方，适合填写英文对照或简短说明。每节副标题\textbf{相互独立}，章节结束后自动清空。
\end{aim*}
\medskip
\begin{aim*}{致谢页用法}{}
在正文末尾调用，格式与封面完全一致：
\smallskip

{\small\texttt{\cmd{makethankspage}}} \quad（默认文字）\\[0.2em]
{\small\texttt{\cmd{makethankspage}[自定义文字]}}
\end{aim*}
\end{column}
\end{columns}
\end{frame}

% ===================== Section 2 =====================
\setsectionsubtitle{Text and Layout}
\section{内容排版}

\begin{frame}{列表与文本框}
\begin{columns}[T]
\begin{column}{0.48\textwidth}
\textbf{无序列表}
\begin{itemize}[left=0pt,label=\textcolor{deepblue}{\ding{113}}]
    \item 使用 \env{itemize} 展示无序条目
    \item 通过 \texttt{enumitem} 宏包定制样式
    \begin{itemize}[left=0pt,label=\textcolor{deepblue}{\ding{237}}]
        \item 支持多级嵌套
        \item 可自定义 \texttt{label} 符号
    \end{itemize}
    \item \textbf{加粗}与\textcolor{tsinghua}{彩色}文字
\end{itemize}
\medskip
\textbf{有序列表}
\begin{enumerate}[left=0pt, label=\arabic*.]
    \item 步骤一：修改封面信息
    \item 步骤二：编写各节内容
    \item 步骤三：XeLaTeX 编译三次
\end{enumerate}
\end{column}
\begin{column}{0.48\textwidth}
\begin{aim*}{浅色文本框 \texttt{aim*}}{}
适合展示研究目标、关键结论等重要信息。支持内嵌列表与公式，背景为浅蓝色，与正文形成层次对比。
\end{aim*}
\end{column}
\end{columns}
\end{frame}

\begin{frame}{分栏布局与三线表}
\begin{columns}[T]
\begin{column}{0.54\textwidth}
\textbf{表格示例（\texttt{booktabs} 三线表）}
\medskip
\begin{table}
\centering
\caption{常用环境速查}
\begin{tabular}{llc}
\toprule
\textbf{环境 / 命令} & \textbf{用途} & \textbf{所需宏包} \\
\midrule
\env{columns}    & 多栏并排布局   & beamer 内置 \\
\env{aim*}       & 浅色高亮文本框 & tcolorbox   \\
\env{lstlisting} & 代码高亮       & listings    \\
\env{algorithm}  & 算法伪代码     & algorithm   \\
\bottomrule
\end{tabular}
\end{table}
\end{column}
\begin{column}{0.42\textwidth}
\begin{aim*}{\texttt{columns} 对齐选项}{}
\begin{itemize}[left=0pt,label=\textcolor{deepblue}{\ding{237}}]
    \item \texttt{[T]} — 顶端对齐（推荐）
    \item \texttt{[c]} — 垂直居中
    \item \texttt{[b]} — 底端对齐
\end{itemize}
\medskip
各栏宽度之和建议不超过 \texttt{0.96\textbackslash textwidth}，过宽会导致内容溢出。
\end{aim*}
\end{column}
\end{columns}
\end{frame}

\begin{frame}{配色说明}
\begin{columns}[T]
\begin{column}{0.50\textwidth}
\begin{aim*}{主题色定义}{}
\begin{tabular}{lll}
\toprule
\textbf{颜色名} & \textbf{RGB 值} & \textbf{示例} \\
\midrule
\texttt{deepblue}  & (51, 70, 139)  & \textcolor{deepblue}{\rule{1.2em}{0.8em}} \\
\texttt{deepred}   & (153, 0, 0)    & \textcolor{deepred}{\rule{1.2em}{0.8em}} \\
\texttt{deepgreen} & (0, 128, 0)    & \textcolor{deepgreen}{\rule{1.2em}{0.8em}} \\
\texttt{halfgray}  & gray 55\%      & \textcolor{halfgray}{\rule{1.2em}{0.8em}} \\
\bottomrule
\end{tabular}
\end{aim*}
\end{column}
\begin{column}{0.46\textwidth}
\begin{aim*}{颜色使用说明}{}
\begin{itemize}[left=0pt,label=\textcolor{deepblue}{\ding{237}}]
    \item \texttt{deepblue} — 主题蓝色，取自\texttt{logo.png} 的主色调 RGB(51,70,139)，用于文本框边框与列表符号
    \item \texttt{deepred} — 深红色，用于强调词及代码自定义高亮
    \item \texttt{deepgreen} — 深绿色，用于代码字符串
    \item \texttt{halfgray} — 浅灰色，用于代码行号
\end{itemize}
\medskip
如需扩展配色，在 \texttt{slide.tex} 头部追加：\\[0.3em]
{\small\texttt{\cmd{definecolor}\{mycolor\}\{RGB\}\{r,g,b\}}}\\[0.2em]
即可在全文通过 \cmd{textcolor}\texttt{\{mycolor\}} 调用。
\end{aim*}
\end{column}
\end{columns}
\end{frame}

% ===================== Section 3 =====================
\setsectionsubtitle{Math and Code}
\section{数学与代码}

\begin{frame}{数学公式}
\begin{columns}[T]
\begin{column}{0.50\textwidth}
\textbf{行内公式：} $E = mc^2$，$\hat{H}\psi = E\psi$
\medskip

\textbf{独立公式：}
\[
\bar{x} = \frac{1}{N}\sum_{i=1}^{N} x_i,\quad
\sigma^2 = \frac{1}{N}\sum_{i=1}^{N}(x_i - \bar{x})^2
\]

\textbf{多行对齐公式：}
\begin{align}
(a+b)^2 &= a^2 + 2ab + b^2 \notag\\
(a-b)^2 &= a^2 - 2ab + b^2 \notag
\end{align}
\end{column}
\begin{column}{0.46\textwidth}
\begin{algorithm}[H]
\caption{梯度下降法}
\begin{algorithmic}[1]
\STATE 初始化参数 $\theta$，学习率 $\alpha$
\WHILE{未收敛}
    \STATE 计算梯度 $g \leftarrow \nabla_\theta L(\theta)$
    \STATE 更新参数 $\theta \leftarrow \theta - \alpha g$
\ENDWHILE
\end{algorithmic}
\end{algorithm}
\end{column}
\end{columns}
\end{frame}

\begin{frame}[fragile]{代码高亮}
\begin{columns}[T]
\begin{column}{0.50\textwidth}
\textbf{Python 示例（\env{lstlisting}）：}
\begin{lstlisting}[language=Python, caption=卡方统计量]
import numpy as np

def chi_square(obs, exp):
    return np.sum(
        (obs - exp)**2 / exp)
\end{lstlisting}
\end{column}
\begin{column}{0.46\textwidth}
\begin{aim*}{lstlisting 常用选项}{}
\begin{itemize}[left=0pt,label=\textcolor{deepblue}{\ding{237}}]
    \item \texttt{language=} — 指定语言\\
    支持 Python、C、TeX 等
    \item \texttt{caption=} — 代码标题
    \item \texttt{numbers=left} — 显示行号
    \item \texttt{basicstyle=} — 字体样式
\end{itemize}
\end{aim*}
\end{column}
\end{columns}
\end{frame}

% ===================== Section 4 =====================
\setsectionsubtitle{Summary}
\section{总结}

\begin{frame}{使用要点速查}
\begin{columns}[T]
\begin{column}{0.48\textwidth}
\begin{aim*}{定制清单}{}
\begin{itemize}[left=0pt,label=\textcolor{deepblue}{\ding{52}}]
    \item 替换 \texttt{logo.png} 为本单位徽标
    \item 替换 \texttt{background1.png} 封面横幅
    \item 替换 \texttt{pic/logo\_ciae.png} 页眉徽标
    \item 修改 \cmd{title}、\cmd{author}、\cmd{institute}
    \item 按需增删 \cmd{section}
    \item 每节前用 \cmd{setsectionsubtitle} 设副标题
\end{itemize}
\end{aim*}
\end{column}
\begin{column}{0.48\textwidth}
\begin{aim*}{编译注意事项}{}
\begin{itemize}[left=0pt,label=\textcolor{deepblue}{\ding{237}}]
    \item 必须使用 \textbf{XeLaTeX} 编译，\textbf{推荐} \texttt{latexmk -xelatex} 自动多次编译
    \item 手动编译时首次或修改结构后需执行\textbf{三次}
    \item 含参考文献时运行顺序：\\[0.3em]
    {\small\texttt{xelatex $\to$ bibtex $\to$ xelatex $\to$ xelatex}}
    \item 编译报错时优先检查特殊字符\\
    是否正确转义（如 \texttt{\textbackslash\%}、\texttt{\textbackslash\$}）
\end{itemize}
\end{aim*}
\end{column}
\end{columns}
\end{frame}

\makethankspage[请各位老师批评指正！]

\end{document}